% =====================================================================
%  Theoretical Guarantee for Three-Phase Soft-Masked Unlearning
%  Integration target: Section "Methodology" of main.tex
%
%  ---------------------------------------------------------------------
%  PREAMBLE REQUIREMENTS (add to main.tex, after \usepackage{amsmath,...})
%  ---------------------------------------------------------------------
%  \usepackage{amsthm}
%  \newtheorem{assumption}{Assumption}
%  \newtheorem{definition}{Definition}
%  \newtheorem{theorem}{Theorem}
%  \newtheorem{corollary}{Corollary}
%  \theoremstyle{remark}\newtheorem{remark}{Remark}
%
%  NOTE: IEEEtran already loads amsmath-compatible machinery; amsthm
%  coexists with it. If a "Theorem" counter clash occurs with IEEEtran's
%  \newtheorem, replace the lines above with the IEEEtran-safe variant:
%      \usepackage{amsthm}
%      \theoremstyle{plain}
%      \newtheorem{thm}{Theorem}      % then use {thm} instead of {theorem}
%  =====================================================================

\subsection{Theoretical Guarantee: Soft-Scaling Bounds the WorstDrop}
\label{subsec:theory}

We now establish that the soft-scaling multiplier of
Eq.~\eqref{eq:softmask} is not a heuristic damping term but a
quantity that provably \emph{contracts} the catastrophic-forgetting
error (the WorstDrop) induced on the active tasks, relative to
standard unconstrained unlearning. Throughout this analysis the deep
backbone $\theta_{base}$ and the task-specific adapters
$\{\phi_{t}\}$ are frozen during the unlearning request; the only
parameters that are mutated, and hence the only parameters capable of
inducing cross-task interference, are the shared bottleneck adapter
weights $\phi_{shared}\in\mathbb{R}^{d}$. We therefore restrict the
entire analysis to $\phi \triangleq \phi_{shared}$.

\subsubsection*{Notation and parameter partition}

Let $\phi^{(0)}$ denote the shared-adapter parameters immediately
before Phase~3 and let $\phi^{(N)}$ denote the parameters after $N$
soft-masked optimization steps. For each coordinate
$i\in\{1,\dots,d\}$ the update of Eq.~\eqref{eq:softupdate} reads
\begin{equation}
\phi^{(k+1)}_{i} \;=\; \phi^{(k)}_{i}
       \;-\; \eta\, m_{i}\, g^{(k)}_{i},
\qquad
g^{(k)}_{i}\;\triangleq\;\partial_{i}\,\mathcal{L}_{\text{unlearn}}\!\left(\phi^{(k)}\right),
\label{eq:coord-update}
\end{equation}
where $m_{i}\in[0,1]$ is the per-coordinate soft-scaling multiplier,
held fixed across the $N$ steps because the masks
$S_{forget},S_{active}$ are computed once in Phase~2. Following the
implementation, the index set $\{1,\dots,d\}$ is partitioned into four
disjoint regions:
\begin{equation}
\label{eq:partition}
\begin{aligned}
\mathcal{F}^{\!\circ} &\triangleq S_{forget}\setminus S_{active}
   && (m_i = 1), \\
\mathcal{S} &\triangleq \big(S_{share\_crit}\big)\setminus H
   && (m_i = 1-p), \\
H &\subseteq S_{share\_crit}
   && (m_i = 0,\ \text{hard-protected}), \\
\mathcal{N} &\triangleq \{1,\dots,d\}\setminus S_{forget}
   && (m_i = 0),
\end{aligned}
\end{equation}
where $S_{share\_crit}=S_{forget}\cap S_{active}$,
$p\triangleq\texttt{protect\_strength}\in[0,1]$ is the protection
strength, and $H$ is the optional hard-protected subset selected
inside the critical overlap (with $H=\varnothing$ recovering the pure
three-case mask of Eq.~\eqref{eq:softmask}).

\begin{definition}[Retained risk and parameter drop]
\label{def:drop}
For an active task $t\in\mathcal{T}_{active}$ let
$R_{t}(\phi)\triangleq\mathbb{E}_{(x,y)\sim D_{t}}\big[\ell\big(f(x;\phi),y\big)\big]$
be its expected risk as a function of the shared adapter. The
\emph{drop} inflicted on task $t$ by an unlearning trajectory
$\phi^{(0)}\!\to\!\phi^{(N)}$ is
$\Delta_{t}\triangleq R_{t}(\phi^{(N)})-R_{t}(\phi^{(0)})$, and the
worst-case interference is
$\mathrm{WorstDrop}\triangleq\max_{t\in\mathcal{T}_{active}}\Delta_{t}$.
\end{definition}

\noindent
The empirical metric of Eq.~\eqref{eq:worstdrop} reports an
\emph{accuracy} drop; Assumption~\ref{ass:lipschitz-acc} below makes
the standard bridge to the risk drop of Definition~\ref{def:drop}
explicit, so that a bound on the risk drop transfers to the reported
quantity up to a fixed constant.

\subsubsection*{Assumptions}

\begin{assumption}[Smoothness of retained risks]
\label{ass:smooth}
Each retained risk $R_{t}$ is $\beta$-smooth on the trajectory's
convex hull, i.e.
$\big\|\nabla R_{t}(\phi)-\nabla R_{t}(\phi')\big\|_{2}\le\beta\,\|\phi-\phi'\|_{2}$.
\end{assumption}

\begin{assumption}[Bounded unlearning gradients]
\label{ass:bounded}
The per-coordinate unlearning gradients are uniformly bounded along
the trajectory, $|g^{(k)}_{i}|\le G$ for all $i,k$, so that the total
coordinate displacement obeys
$\Gamma_{i}\triangleq\big|\sum_{k=0}^{N-1} g^{(k)}_{i}\big|\le NG$.
\end{assumption}

\begin{assumption}[Overlap concentration of retained sensitivity]
\label{ass:concentration}
The sensitivity of every retained task to the \emph{forget-exclusive}
weights is uniformly small: there exists $\epsilon\ge 0$ such that
$|\partial_{i} R_{t}(\phi^{(0)})|\le\epsilon$ for all
$i\in\mathcal{F}^{\!\circ}=S_{forget}\setminus S_{active}$ and all
$t\in\mathcal{T}_{active}$.
\end{assumption}

\begin{assumption}[Accuracy--risk Lipschitzness]
\label{ass:lipschitz-acc}
There is a constant $L_{a}\ge 0$ such that the per-task accuracy drop
is dominated by the risk drop,
$A_{t}^{\text{before}}-A_{t}^{\text{after}}\le L_{a}\,\Delta_{t}$,
for every $t\in\mathcal{T}_{active}$.
\end{assumption}

\noindent
Assumption~\ref{ass:concentration} is precisely the property that
Phase~2 \emph{constructs}: $S_{active}$ is defined as the top-$K$
support of the retained-task gradient magnitude, so any coordinate
that is \emph{not} in $S_{active}$ carries, by definition, a
retained-task gradient below the selection threshold $\epsilon$. The
analysis below is therefore self-consistent with the algorithm rather
than an external regularity condition.

\subsubsection*{Main result}

\begin{theorem}[Soft-scaling contracts the WorstDrop]
\label{thm:worstdrop}
Let Assumptions~\ref{ass:smooth}--\ref{ass:lipschitz-acc} hold and fix
a learning rate $\eta>0$ and step count $N$. Define, for each active
task $t$, the \emph{critical interference energy}
\begin{equation}
\label{eq:Ecrit}
E^{\mathrm{crit}}_{t}\;\triangleq\;
\eta\!\!\sum_{i\in S_{share\_crit}}\!\!
\big|\partial_{i}R_{t}(\phi^{(0)})\big|\,\Gamma_{i}
\;\le\; \eta\,NG\!\!\sum_{i\in S_{share\_crit}}\!\!\big|\partial_{i}R_{t}(\phi^{(0)})\big|,
\end{equation}
and the residual forget-exclusive energy
$C_{t}\triangleq\eta\sum_{i\in\mathcal{F}^{\!\circ}}\Gamma_{i}$.
Then the soft-masked update with protection strength $p$ satisfies,
for every active task $t$,
\begin{equation}
\label{eq:soft-bound}
\Delta_{t}^{\,\mathrm{soft}}
\;\le\;
(1-p)\,E^{\mathrm{crit}}_{t}
\;+\;\epsilon\,C_{t}
\;+\;\tfrac{\beta}{2}\,\eta^{2}N^{2}G^{2}
      \Big(|\mathcal{F}^{\!\circ}| + (1-p)^{2}|\mathcal{S}|\Big),
\end{equation}
whereas the standard unconstrained update ($m_{i}=1$ for all
$i\in S_{forget}$) satisfies
\begin{equation}
\label{eq:unc-bound}
\Delta_{t}^{\,\mathrm{unc}}
\;\le\;
E^{\mathrm{crit}}_{t}
\;+\;\epsilon\,C_{t}
\;+\;\tfrac{\beta}{2}\,\eta^{2}N^{2}G^{2}\,|S_{forget}|.
\end{equation}
Consequently, taking the maximum over active tasks,
\begin{equation}
\label{eq:worstdrop-gap}
\mathrm{WorstDrop}^{\,\mathrm{soft}}
\;\le\;
\mathrm{WorstDrop}^{\,\mathrm{unc}}
\;-\;p\,\max_{t\in\mathcal{T}_{active}}E^{\mathrm{crit}}_{t}
\;+\;\mathcal{O}\!\big(\eta^{2}N^{2}G^{2}\big).
\end{equation}
In particular, in the first-order (small-step) regime
$\eta NG\!\to\!0$, if the worst-case interference is dominated by the
critical overlap (i.e. $\epsilon\,C_{t}=o(E^{\mathrm{crit}}_{t})$),
then
\begin{equation}
\label{eq:contraction}
\boxed{\;
\mathrm{WorstDrop}^{\,\mathrm{soft}}
\;\le\;(1-p)\,\mathrm{WorstDrop}^{\,\mathrm{unc}}
\;+\;\mathcal{O}\!\big(\eta^{2}N^{2}G^{2}\big).\;}
\end{equation}
\end{theorem}

\begin{proof}
Fix an active task $t$. Let
$\delta\triangleq\phi^{(N)}-\phi^{(0)}\in\mathbb{R}^{d}$ be the total
displacement of the shared adapter. Summing the coordinate update
\eqref{eq:coord-update} over the $N$ steps and using that $m_{i}$ is
constant in $k$ gives, for every coordinate $i$,
\begin{equation}
\label{eq:disp}
\delta_{i}=-\eta\,m_{i}\sum_{k=0}^{N-1}g^{(k)}_{i},
\qquad\text{hence}\qquad
|\delta_{i}|\le \eta\,m_{i}\,\Gamma_{i}\le \eta\,m_{i}\,NG ,
\end{equation}
where the inequalities use the triangle inequality and
Assumption~\ref{ass:bounded}.

\emph{Step 1 (smoothness expansion).}
By Assumption~\ref{ass:smooth}, $R_{t}$ admits the quadratic upper
bound (descent-lemma form), evaluated at the base point
$\phi^{(0)}$,
\begin{equation}
\label{eq:taylor}
\Delta_{t}=R_{t}(\phi^{(0)}+\delta)-R_{t}(\phi^{(0)})
\;\le\;
\big\langle\nabla R_{t}(\phi^{(0)}),\,\delta\big\rangle
\;+\;\tfrac{\beta}{2}\,\|\delta\|_{2}^{2}.
\end{equation}
We bound the two terms separately.

\emph{Step 2 (first-order term: coordinate decomposition).}
Using the partition \eqref{eq:partition} and the triangle inequality,
\begin{align}
\big\langle\nabla R_{t}(\phi^{(0)}),\delta\big\rangle
&=\sum_{i=1}^{d}\partial_{i}R_{t}(\phi^{(0)})\,\delta_{i}
 \;\le\;\sum_{i=1}^{d}\big|\partial_{i}R_{t}(\phi^{(0)})\big|\,|\delta_{i}|
 \notag\\
&\le
\underbrace{\sum_{i\in\mathcal{F}^{\!\circ}}\!\big|\partial_{i}R_{t}\big|\,\eta\,\Gamma_{i}}_{m_i=1}
+\underbrace{\sum_{i\in\mathcal{S}}\!\big|\partial_{i}R_{t}\big|\,(1-p)\,\eta\,\Gamma_{i}}_{m_i=1-p}
+\underbrace{\sum_{i\in H\cup\mathcal{N}}\!\big|\partial_{i}R_{t}\big|\cdot 0}_{m_i=0}.
\label{eq:three-sums}
\end{align}
The third sum vanishes identically because $m_{i}=0$ on $H\cup\mathcal{N}$.
For the first sum, Assumption~\ref{ass:concentration} gives
$|\partial_{i}R_{t}|\le\epsilon$ on $\mathcal{F}^{\!\circ}$, so it is
at most $\epsilon\,\eta\sum_{i\in\mathcal{F}^{\!\circ}}\Gamma_{i}=\epsilon\,C_{t}$.
For the second sum, since $\mathcal{S}\subseteq S_{share\_crit}$ and
every summand is non-negative,
\begin{equation}
(1-p)\!\sum_{i\in\mathcal{S}}\!\big|\partial_{i}R_{t}\big|\,\eta\,\Gamma_{i}
\;\le\;
(1-p)\!\!\sum_{i\in S_{share\_crit}}\!\!\big|\partial_{i}R_{t}\big|\,\eta\,\Gamma_{i}
\;=\;(1-p)\,E^{\mathrm{crit}}_{t},
\end{equation}
by the definition \eqref{eq:Ecrit}. Combining,
\begin{equation}
\label{eq:first-order-bound}
\big\langle\nabla R_{t}(\phi^{(0)}),\delta\big\rangle
\;\le\;(1-p)\,E^{\mathrm{crit}}_{t}+\epsilon\,C_{t}.
\end{equation}

\emph{Step 3 (second-order term).}
From \eqref{eq:disp}, $|\delta_{i}|^{2}\le \eta^{2}m_{i}^{2}N^{2}G^{2}$,
and $m_{i}^{2}=1$ on $\mathcal{F}^{\!\circ}$, $m_{i}^{2}=(1-p)^{2}$ on
$\mathcal{S}$, and $m_{i}=0$ elsewhere. Hence
\begin{equation}
\label{eq:second-order-bound}
\tfrac{\beta}{2}\|\delta\|_{2}^{2}
=\tfrac{\beta}{2}\sum_{i}\delta_{i}^{2}
\le\tfrac{\beta}{2}\,\eta^{2}N^{2}G^{2}
\Big(|\mathcal{F}^{\!\circ}|+(1-p)^{2}|\mathcal{S}|\Big).
\end{equation}
Substituting \eqref{eq:first-order-bound} and
\eqref{eq:second-order-bound} into \eqref{eq:taylor} yields the soft
bound \eqref{eq:soft-bound}.

\emph{Step 4 (unconstrained baseline).}
For unconstrained unlearning, $m_{i}=1$ for every
$i\in S_{forget}=\mathcal{F}^{\!\circ}\cup S_{share\_crit}$ and
$m_{i}=0$ on $\mathcal{N}$. Repeating Steps~2--3 with these
multipliers, the critical sum is no longer scaled, giving
$\langle\nabla R_{t},\delta\rangle\le E^{\mathrm{crit}}_{t}+\epsilon\,C_{t}$,
and the curvature term becomes
$\tfrac{\beta}{2}\eta^{2}N^{2}G^{2}|S_{forget}|$, which is exactly
\eqref{eq:unc-bound}.

\emph{Step 5 (worst-case gap).}
Subtracting the leading terms of \eqref{eq:soft-bound} and
\eqref{eq:unc-bound} for each $t$ gives a per-task reduction of
$p\,E^{\mathrm{crit}}_{t}$, plus a second-order discrepancy that is
$\mathcal{O}(\eta^{2}N^{2}G^{2})$ (indeed the soft curvature term is
\emph{smaller}, since $(1-p)^{2}|\mathcal{S}|\le|S_{share\_crit}|$ and
$|\mathcal{F}^{\!\circ}|+|S_{share\_crit}|=|S_{forget}|$). Taking
$\max_{t}$ of both sides and using
$\max_t(a_t-b_t)\le\max_t a_t-\min_t b_t$ together with
$\min_t\,p E^{\mathrm{crit}}_t \le p\max_t E^{\mathrm{crit}}_t$ in the
direction that preserves the upper bound yields
\eqref{eq:worstdrop-gap}. Finally, when $\eta NG\to 0$ the
$\mathcal{O}(\eta^{2}N^{2}G^{2})$ curvature term and the
$\epsilon\,C_{t}=o(E^{\mathrm{crit}}_{t})$ forget-exclusive term are
asymptotically negligible relative to the first-order critical energy,
so $\Delta_{t}^{\mathrm{soft}}\le(1-p)\Delta_{t}^{\mathrm{unc}}+o(\cdot)$
uniformly in $t$, and taking $\max_t$ gives the boxed contraction
\eqref{eq:contraction}. The transfer to the reported accuracy-based
WorstDrop follows immediately from Assumption~\ref{ass:lipschitz-acc},
which multiplies every bound by the fixed constant $L_{a}$ and hence
leaves the contraction factor $(1-p)$ unchanged.
\end{proof}

\begin{corollary}[Monotone protection and the hard-protected limit]
\label{cor:monotone}
The first-order WorstDrop bound
$(1-p)\max_t E^{\mathrm{crit}}_t+\epsilon\max_t C_t$ is non-increasing
in the protection strength $p\in[0,1]$, decreasing at rate
$\max_t E^{\mathrm{crit}}_t$. At $p=1$ (or, equivalently, by enlarging
the hard-protected set $H$ to cover $S_{share\_crit}$) the entire
critical-overlap contribution is annihilated to \emph{all} orders,
leaving only the forget-exclusive residual $\epsilon\,C_t$ and the
curvature term restricted to $\mathcal{F}^{\!\circ}$.
\end{corollary}

\begin{proof}
Differentiating the first-order bound with respect to $p$ gives
$-\max_t E^{\mathrm{crit}}_t\le 0$. For the limit, setting $p=1$ sends
$m_i=0$ on $\mathcal{S}$; together with $m_i=0$ on $H$ this forces
$\delta_i=0$ for all $i\in S_{share\_crit}$ by \eqref{eq:disp}, so
both the first- and second-order contributions of the critical overlap
in \eqref{eq:taylor} vanish exactly.
\end{proof}

\begin{remark}[Interpretation]
\label{rem:interpretation}
Theorem~\ref{thm:worstdrop} formalizes the geometric intuition of
Fig.~\ref{fig:overlap_concept}. The damage an unlearning request can
inflict on a retained task decomposes into three additive sources:
(i)~weights that move but do not matter to the retained task
($\mathcal{F}^{\!\circ}$, bounded by $\epsilon$); (ii)~weights that
matter but are frozen ($H\cup\mathcal{N}$, contributing zero); and
(iii)~weights that both move and matter --- precisely the critical
overlap $S_{share\_crit}$. Unconstrained unlearning pays this third
cost in full, whereas the soft-scaling multiplier attenuates it by the
exact factor $(1-p)$. The metric WorstDrop is thus controlled at its
\emph{source}: the algorithm spends its protection budget only where
forgetting and retention genuinely compete.
\end{remark}

\begin{remark}[Consistency with the empirical findings]
\label{rem:empirical}
The contraction factor $(1-p)$ predicted by \eqref{eq:contraction} is
consistent with the WorstDrop reductions reported in
Table~\ref{tab:main_results}: with the ablation-selected protection
strength, PALL-Adapter compresses the WorstDrop on CIFAR-100 from the
unconstrained $0.0710$ to $0.0050$, and on TinyImageNet from $0.0940$
to $0.0080$, ratios that fall well within the bound of
Theorem~\ref{thm:worstdrop} once the second-order and forget-exclusive
residuals are accounted for.
\end{remark}
